% Generated by scripts/assemble.py; corrections belong in source/corrections.json.
\examhead{94}{河海大学 2025 年数学分析试卷}{93}
% p094-b01; source PDF page 94
\par\Needspace{5\baselineskip}
\blockanchor{p094-b01}
\majorquestion{一}判断题. 判断下列命题是否正确并说明理由. 若判断某命题错误, 则务必给出反例或说明理由, 否则该题仅得 2 分.\par
% p094-b02; source PDF page 94
\blockanchor{p094-b02}
\question{1}(5 分) 数列 $\{a_{n}\}$ 为实数列, 如果对所有的 $n, p \in \mathbb{N}^{+}$, 都有 $|a_{n+p} - a_{n}| < \frac{p}{n}$, 那么 $\{a_{n}\}$ 为基本数列(Cauchy 列).\par
% p094-b03; source PDF page 94
\blockanchor{p094-b03}
\question{2}(5 分) 如果连续函数 $f(x)$ 在 $(a, b)$ 上没有零点, 那么 $f(x)$ 在 $(a, b)$ 内不变号.\par
% p094-b04; source PDF page 94
\blockanchor{p094-b04}
\question{3}(5 分) 如果函数 $f(x)$ 在 $[a, b]$ 上可积, 那么 $f(x)$ 在 $[a, b]$ 上至多只有有限个间断点.\par
% p094-b05; source PDF page 94
\blockanchor{p094-b05}
\question{4}(5 分) 如果级数 $\sum_{n=1}^{\infty} a_n$ 收敛, 实数列 $\{b_n\}$ 满足 $\lim_{n \to \infty} \frac{b_n}{a_n} = 0$, 那么 $\sum_{n=1}^{\infty} b_n$ 也收敛.\par
% p094-b06; source PDF page 94
\par\Needspace{5\baselineskip}
\blockanchor{p094-b06}
\majorquestion{二}叙述题.\par
% p094-b07; source PDF page 94
\blockanchor{p094-b07}
\question{5}(5 分) 写出命题“当 $x \to x_0^-$ 时, 函数 $f(x)$ 的极限为正无穷大”的否定命题的分析表述.\par
% p094-b08; source PDF page 94
\blockanchor{p094-b08}
\question{6}(5 分) 叙述函数 $f(x)$ 在区间 $I$ 上一致连续的定义.\par
% p094-b09; source PDF page 94
\blockanchor{p094-b09}
\question{7}(5 分) 叙述函数 $f(x)$ 在 $[a, b]$ 上 Riemann 可积的定义.\par
% p094-b10; source PDF page 94
\blockanchor{p094-b10}
\question{8}(5 分) 叙述判定函数项级数 $\sum_{n=1}^{\infty} a_n(x) b_n(x)$ 在 $D$ 上一致收敛的 Dirichlet 判别法.\par
% p094-b11; source PDF page 94
\par\Needspace{5\baselineskip}
\blockanchor{p094-b11}
\majorquestion{三}解答与证明题.\par
% p094-b12; source PDF page 94
\blockanchor{p094-b12}
\question{9}(10 分) 求极限 $\lim_{x\to 0}\frac{\frac{x^2}{2} - \sqrt{1 + x^2} + 1}{(\cos x - e^{x^2})\tan^2x}.$\par
% p094-b13; source PDF page 94
\blockanchor{p094-b13}
\question{10}(10 分) 求不定积分 $\int \frac{x^2 \arcsin x}{\sqrt{1 - x^2}} \, \mathrm{d}x$.\par
% p094-b14; source PDF page 94
\par\Needspace{6\baselineskip}
\blockanchor{p094-b14}
\question{11}(10 分) 设 $p > 0$, 讨论级数\par
\nopagebreak[4]
% p094-b15; source PDF page 94
\blockanchor{p094-b15}
\displayformula{\sum_ {n = 1} ^ {\infty} \ln \left(1 + \frac {(- 1) ^ {n - 1}}{n ^ {p}}\right)}
% p094-b16; source PDF page 94
\blockanchor{p094-b16}
是否收敛？如果收敛的话，是条件收敛还是绝对收敛？\par
% p094-b17; source PDF page 94
\par\Needspace{6\baselineskip}
\blockanchor{p094-b17}
\question{12}(10 分) 函数 $z = z(x, y)$ 具有二阶连续偏导数, 且满足方程\par
\nopagebreak[4]
% p094-b18; source PDF page 94
\blockanchor{p094-b18}
\displayformula{\frac {\partial^ {2} z}{\partial x ^ {2}} - 2 \frac {\partial^ {2} z}{\partial x \partial y} + \frac {\partial^ {2} z}{\partial y ^ {2}} = 0.}
% p094-b19; source PDF page 94
\blockanchor{p094-b19}
现做自变量变换 $u = x - y, v = x + y$ 和因变量变换 $w = xy - z$, 求 $w$ 作为 $u$ 和 $v$ 的函数所满足的方程.\par
% p094-b20; source PDF page 94
\par\Needspace{6\baselineskip}
\blockanchor{p094-b20}
\question{13}(10 分) 设 $h > 0$, 求\par
\nopagebreak[4]
% p094-b21; source PDF page 94
\blockanchor{p094-b21}
\displayformula{\iint_ {\Sigma} x ^ {3} \mathrm{d} y \mathrm{d} z + y ^ {3} \mathrm{d} z \mathrm{d} x + z ^ {3} \mathrm{d} x \mathrm{d} y.}
% p094-b22; source PDF page 94
\blockanchor{p094-b22}
其中 $\Sigma$ 是锥面 $z = \sqrt{x^2 + y^2}$ 在 $0 \leq z \leq h$ 的部分, 方向取上侧.\par
% p094-b23; source PDF page 94
\blockanchor{p094-b23}
\question{14}(15 分) 函数 $f(x)$ 在 $[0,1]$ 上二阶连续可导, 且 $|f(x)| \leq \min \{1 - \cos x, \sin \pi x\}$, $\forall x \in [0,1]$. 证明:\par
% p094-b24; source PDF page 94
\blockanchor{p094-b24}
(1) 存在 $\xi \in (0,1)$, 使得 $f'(\xi) = 0$.\par
% p094-b25; source PDF page 94
\blockanchor{p094-b25}
(2) 存在 $\eta \in (0,1)$, 使得 $f''(\eta) = 0$.\par
% p094-b26; source PDF page 94
\blockanchor{p094-b26}
\question{15}(15 分) 设 $f(x, y) = \frac{1}{(1 + xy)(1 + x)}$, $(x, y) \in [0, +\infty) \times [0, +\infty)$.\par
% p094-b27; source PDF page 94
\blockanchor{p094-b27}
(1) 证明 $\int_0^{+\infty} f(x, y) \mathrm{d}x$ 关于 $y \in (0, +\infty)$ 上是内闭一致收敛的.\par
% p094-b28; source PDF page 94
\blockanchor{p094-b28}
(2) 设 $g(y) = \int_{0}^{+\infty} \frac{\ln(1 + xy)}{x(1 + x)} \, \mathrm{d}x, y \in (0, +\infty)$, 求 $g'(y)$ 及 $g(1)$.\par
% p095-b00; source PDF page 95
\blockanchor{p095-b00}
\question{16}(15 分) 解答如下问题:\par
% p095-b01; source PDF page 95
\blockanchor{p095-b01}
(1) 求幂级数 $\sum_{n=1}^{\infty} \frac{x^n}{n}$ 的收敛域以及和函数.\par
% p095-b02; source PDF page 95
\par\Needspace{6\baselineskip}
\blockanchor{p095-b02}
(2) $f(x) = \sum_{n=1}^{\infty} \frac{x^n}{n^2}, x \in [-1, 1]$. 证明: 对任意的 $x \in (0, 1)$, 都有\par
\nopagebreak[4]
% p095-b03; source PDF page 95
\blockanchor{p095-b03}
\displayformula{f (x) + f (1 - x) + \ln x \cdot \ln (1 - x) = \frac {\pi^ {2}}{6}.}
% p095-b04; source PDF page 95
\par\Needspace{6\baselineskip}
\blockanchor{p095-b04}
\question{17}(15 分) 设数列 $\{a_{n}\}$ 满足 $\lim_{n\to \infty}a_n = a,$ 数列 $\{p_n\}$ 满足 $p_k > 0, k = 1,2,\dots$，且\par
\nopagebreak[4]
% p095-b05; source PDF page 95
\blockanchor{p095-b05}
\displayformula{\lim _ {n \rightarrow \infty} \frac {p _ {n}}{p _ {1} + p _ {2} + \cdots + p _ {n}} = 0.}
% p095-b06; source PDF page 95
\par\Needspace{6\baselineskip}
\blockanchor{p095-b06}
证明：\par
\nopagebreak[4]
% p095-b07; source PDF page 95
\blockanchor{p095-b07}
\displayformula{\lim _ {n \rightarrow \infty} \frac {p _ {1} a _ {n} + p _ {2} a _ {n - 1} + \cdots + p _ {n} a _ {1}}{p _ {1} + p _ {2} + \cdots + p _ {n}} = a.}
